\begin{table}[tbp] \centering
\newcolumntype{R}{>{\raggedleft\arraybackslash}X}
\newcolumntype{L}{>{\raggedright\arraybackslash}X}
\newcolumntype{C}{>{\centering\arraybackslash}X}

\caption{French--German turnout gap collapse on the absinthe ban (Vote \#68)}
\label{tab:T6_turnout_gap}
\begin{tabularx}{\linewidth}{@{}lCCC@{}}

\toprule
{ }&{French (N=5)}&{German (N=20)}&{Gap (Fr--Ge)} \tabularnewline
\midrule \addlinespace[\belowrulesep]
Baseline turnout (median across 14 placebo votes 1907--1910)&44.40&56.62&--12.22 \tabularnewline
Vote \#68 turnout (5 July 1908, absinthe ban)&47.92&48.02&--0.11 \tabularnewline
Change (gap collapse, row 2 minus row 1)&+3.52&--8.60&\textbf{+12.12} \tabularnewline
\bottomrule \addlinespace[\belowrulesep]

\end{tabularx}
\\ \parbox{\linewidth}{\footnotesize French = cantons with French language share >50\% of (German+French) speakers (NE, GE, VD, FR, VS; N=5). German = remaining 20 cantons. Baseline turnout = canton-level median across 14 federal referenda 1907--1910 (placebo set used to detect mobilization deviation on \#68). Gap entries are arithmetic differences between French and German group means.}
\end{table}
